%% =====================================================================
%%  T-03-02  The Operator Profile
%%  -------------------------------------------------------------------
%%  One idea per file. Core is mandatory. Slots in canonical order.
%%  Max 700 words. No layout commands. No filler. New source material
%%  for this entry goes in sources/<BOOK>/T-03-02.tex -- never by
%%  rewriting this file (docs/RULES.md R0.1).
%% =====================================================================
%% @kind: topic
%% @id: T-03-02
%% @title: The Operator Profile
%% @subtitle: What the person working you usually looks like
%% @chapter: c03
%% @order: 20
%% @status: stable
%% @sources: B1
%% @updated: 2026-09-19

\topic{T-03-02}{The Operator Profile}{What the person working you usually looks like}

\begin{core}
The experienced operator does not present as dangerous. They present as warm, unlucky, generous with gratitude and slightly in need --- a combination that disarms exactly the faculties required to notice them.
\end{core}
\begin{mechanism}
Each feature of the profile does specific work. Warmth and early intimacy accelerate the relationship past the point where ordinary suspicion operates, and they create an obligation to reciprocate. A story of former wealth explains present need without implying incompetence, which makes the request for money legible and sympathetic. Demonstrative gratitude pre-pays the emotional reward so that the target feels the exchange is already balanced. Insistence on their own trustworthiness substitutes a claim for evidence. The profile works because every element is something a decent person also displays; the difference is the concentration and the timing.
\end{mechanism}
\begin{tells}
  \item Intimacy arrives faster than the relationship can support.
  \item They were once wealthy, or once important, and something was done to them.
  \item Gratitude is expressed at length and in advance, before anything has been given.
  \item Trustworthiness is asserted repeatedly and unprompted.
  \item The benefit you stand to gain is restated several times by the person proposing the arrangement.
  \item Every story has them as the person who was wronged.
\end{tells}
\begin{moves}
  \item Watch for the profile rather than for a single tell --- no individual feature is proof, and the pattern is.
  \item Note where each feature appears in the sequence: warmth first, story second, gratitude third, request fourth.
  \item Treat any request involving money as the point at which the profile is confirmed or dismissed.
  \item Do not test them with a confrontation. Test them with a delay: the operator's response to being asked to wait is more informative than their response to being accused.
\end{moves}
\begin{counters}
  \item Slow the relationship to its natural rate. Intimacy that is refused acceleration usually reveals itself.
  \item Require the claim of trustworthiness to be evidenced rather than stated --- references, a written agreement, a verifiable history.
  \item Separate gratitude from obligation. Thanks is not a payment and creates no debt.
  \item Keep money and personal warmth in different accounts, and notice anyone who tries to move between them quickly.
\end{counters}
\begin{limits}
The profile describes the small-scale operator: the borrower, the deal-maker, the plausible stranger. It is much less useful against institutional or long-relationship manipulation, where the operator has years of genuine reliability behind them and none of the surface markers. Used indiscriminately it also produces a person who cannot accept kindness, which is a real cost --- warmth, misfortune and gratitude are all common in honest people, and the base rate of the profile is low.
\end{limits}

\sources{B1}
\seesources
\seealso{T-03-01, T-06-01, T-25-01}
